\documentclass[a4paper,10pt]{report}\input{../0.Préambule/Préambule_cours_prof}\begin{document}

\pagestyle{empty}
\newpage
\noindent NOM, Prénom et classe : \dotfill

\section*{Les proportions et les pourcentages}
\addcontentsline{toc}{section}{Les proportions et les pourcentages}

\noindent \textit{L'usage de la calculatrice est \textbf{interdite}.} \hfill \textit{Durée : 55 minutes}

\paragraph{Exercice 1}: \hfill \textbf{6 points}

\medskip
\noindent Le père de Lucas achète une tranche de $150$g de Roquefort Bio à $3$ euros,
\begin{enumerate}
\item Combien coûtent $100$g de ce fromage ?
\item Combien coûtent $300$g de ce fromage ?
\item Quel est le prix au kilogramme de ce fromage ?
\end{enumerate}

\paragraph{Exercice 2}: \hfill \textbf{5 points}

\medskip
\noindent Gabriella participe à une épreuve de course à pieds. Elle court à la même vitesse tout au long du parcours. Les résultats des temps de passage figurent dans la tableau suivant :

\medskip
\noindent \begin{minipage}{9cm}
\begin{enumerate}
\item Ce tableau est-il un tableau de proportionnalité ? Justifie
\item Si oui, calcule le coefficient de proportionnalité ?
\item Quelle distance Gabriella parcourt-elle en une heure ?
\end{enumerate}
\end{minipage}
\hspace{2mm}
\begin{minipage}{7.5cm}
\begin{center}
\renewcommand{\arraystretch}{1.5}
\begin{tabularx}{\linewidth}{|p{4cm}|*{4}{>{\centering \arraybackslash}X|}}
\hline
\cellcolor{lightgray!50}{
Kilomètres effectués} & $4,5$ & $6$ & $11$ & $14$\\\hline
\cellcolor{lightgray!50}{Temps en minutes} & $27$ & $36$ & $66$ & $84$\\\hline
\end{tabularx}
\end{center}
\end{minipage}

\paragraph{Exercice 3}: \hfill \textbf{5 points}

\medskip
\noindent Calcule les pourcentages suivants :
\begin{enumerate}
\begin{tabular}{m{5cm}m{5cm}m{5cm}}
\item $25\%$ de $12$ &
\item $10\%$ de $1420$g &
\item $75\%$ de $52$L \\
\item $50\%$ de $438$m &
\item $60\%$ de $200$km &\\
\end{tabular}
\end{enumerate}

\paragraph{Exercice 4}: \hfill \textbf{4 points}

\medskip
\noindent Au mois de janvier, le loyer des parents de François a augmenté de $4 \%$. Le loyer était de $800$ euros.
\begin{enumerate}
\item A combien s'élève l'augmentation ?
\item Quel est alors le montant du nouveau loyer ?
\end{enumerate}

\paragraph{Exercices bonus}: \hfill \textbf{0,5 point}

\medskip
Lola s'amuse à plier une feuille de papier. Sachant que la feuille mesure $0,1$mm d'épaisseur, quelle épaisseur devrait trouver Lola si elle pliait la feuille sur elle-même $20$ fois de suite ?

 \hfill \textbf{0,5 point}

Un fermier a des poules et des lapins. En regardant tous les animaux, il voit 5 têtes et 16 pattes. Combien le fermier a-t-il de lapins et de poules ?

\hfill \textbf{0,5 point}

Pendule, ma pendule, il est six heures et demie du soir.
Quel est donc l'angle que forment tes aiguilles ?

\hfill \textbf{0,5 point}

On dispose de 7 petites boîtes identiques. L'une n'a pas le même poids que les autres. En utilisant une balance à plateaux, décrire les pesées à effectuer pour retrouver cette boîte en un nombre minimum de pesées.

\end{document}